% ===========================================================================
\ReportSection[sec:concept]{Proposed System Concept}

\ReportSubsection{Two-computer architecture}

\FigureTitle[fig:arch]{Logical architecture: Computer~1 (NAS and private cloud) and Computer~2 (ML nodes)}
\begin{center}
\begin{tikzpicture}[
  svc/.style={draw, rounded corners=1.5pt, font=\scriptsize, align=center,
              minimum height=0.8cm, minimum width=3.4cm, inner sep=2pt},
  host/.style={draw, rounded corners=3pt, inner sep=7pt},
  lbl/.style={font=\small\itshape},
  arr/.style={-{Stealth[length=2mm]}, thick},
  darr/.style={-{Stealth[length=2mm]}, thick, dashed},
  bi/.style={{Stealth[length=2mm]}-{Stealth[length=2mm]}, thick}
]
\node[svc, minimum width=2.4cm] (users) at (0,-1) {Users:\\browser, IDE,\\scripts};
\node[svc] (gw)  at (5.2,-1) {LLM gateway, chat UI,\\retrieval service};
\node[svc] (vdb) at (5.2,-2.2) {Vector database,\\model and artifact registry};
\node[svc] (fs)  at (5.2,-3.4) {File share (SMB/NFS),\\snapshots, backups};
\node[svc] (lic) at (5.2,-4.6) {License servers,\\private-cloud VMs};
\begin{scope}[on background layer]
  \node[host, fit=(gw)(vdb)(fs)(lic), label={[lbl]above:Computer 1: NAS / private cloud}] (c1) {};
\end{scope}
\node[svc, minimum width=4.2cm, minimum height=1.3cm] (na) at (12.7,-1) {ML Node A (primary)\\serving: inference,\\embeddings, reranker};
\node[svc, minimum width=4.2cm, minimum height=1.3cm] (nb) at (12.7,-3.4) {ML Node B (secondary)\\fine-tuning, experiments;\\standby for Node A};
\node[lbl, above=0.15cm of na] {Computer 2: ML nodes};
\draw[arr] (users) -- (gw);
\draw[arr] (gw.east) -- node[above, font=\scriptsize]{inference API} (na.west);
\draw[darr] (gw.east) ++(0.5,0) |- node[pos=0.25, right, font=\scriptsize]{failover} ([yshift=0.35cm]nb.west);
\draw[bi] (fs.east) -- node[below, font=\scriptsize]{datasets, checkpoints, backups} (nb.west |- fs.east);
\end{tikzpicture}
\end{center}
{\footnotesize Every node mounts the NAS over the storage network and keeps its own
local copy of the approved models, so serving continues if Computer~1 is
briefly unavailable. A second gateway instance can run on Node~B (Section~\ref{sec:resil}).}

\ReportSubsection{Computer 1: roles}
\begin{itemize}
  \item \textbf{File and data services:} SMB share with Active Directory or
    Kerberos authentication, NFS exports to the ML nodes, snapshot schedule,
    quotas and an encrypted backup target.
  \item \textbf{License servers:} hosted in virtual machines or on the host.
    \emph{Check each vendor's terms}: some license managers are bound to a
    hardware or MAC address and may restrict virtualization.
  \item \textbf{Private cloud:} a hypervisor (for example Proxmox or KVM on a
    FIPS-capable Linux) running the specific software \PH{list}, each in its
    own virtual machine or container, with snapshots and backup.
  \item \textbf{LLM application tier:} gateway, chat interface, retrieval
    service, vector database, model and artifact registry, observability.
\end{itemize}

\ReportSubsection[sec:licenses]{Hosting Autodesk and Bluebeam}
\TodoAnchor{aoa-licenses}

The stated software to host is Autodesk products and Bluebeam. These need to be checked
with the vendors or reseller before the Computer~1 design is frozen; the points below are
the questions to ask, not confirmed facts.

\TableTitle{Autodesk and Bluebeam: what to verify}
{\footnotesize
\begin{xltabular}{\textwidth}{@{}Y{0.636} Y{1.364} Y{1.0}@{}}
  \toprule
  Topic & Question & Why it matters \\
  \midrule
  \endhead
  Licensing model & Are the licenses named-user subscriptions that sign in through the vendor's cloud, or network (floating) licenses served from a local license manager? (Autodesk and Bluebeam have largely moved to named-user subscriptions; legacy agreements may still use a local license manager.) & If named-user, a local license server may not be needed; if network, the license manager runs on Computer~1. \\
  Internet sign-in & Does the product need periodic internet check-in or sign-in? & An enclave or air-gapped network may not allow it. \\
  Data and collaboration services & Which services are needed (for example Autodesk Vault for data management, Bluebeam collaboration)? Is a self-hosted option available and authorized for CUI? & Cloud versions sit outside the CUI boundary unless authorized. \\
  Server platform & Do the server components require Windows Server and SQL Server? & Adds Windows Server, database and client-access licensing, and may favor a Hyper-V host. \\
  Virtualization and hardware binding & Does the vendor support running the license manager in a virtual machine, and is it bound to a host ID or MAC address? & May require a fixed virtual MAC or a dedicated host. \\
  \bottomrule
\end{xltabular}}

\ReportSubsection{Computer 2: ML nodes and workloads}

\TableTitle{How each node is used}
{\small
\begin{tabularx}{\textwidth}{@{}Y{0.778}Y{1.111}Y{1.111}@{}}
  \toprule
  & Single-node configs & Dual-node configs \\
  \midrule
  Serving & Node serves all users. & Node~A serves; Node~B is warm standby. \\
  Fine-tuning & Time-shared with serving (off-hours) or on a separate GPU if the node has two or more. & Runs on Node~B concurrently with serving on Node~A. \\
  Experiments and evaluation & Same as fine-tuning. & Node~B; does not disturb production latency. \\
  Failure of the node & Service down until repaired or the cloud fallback is used. & Gateway routes to Node~B; training pauses until Node~A returns. \\
  Maintenance (patching, driver or model upgrades) & Planned downtime. & Rolling: patch Node~B first, then switch and patch Node~A. \\
  \bottomrule
\end{tabularx}}

A dual-node design is worth its extra cost mainly for three reasons: (1)
training cannot degrade user-facing latency, (2) there is a failover path, and
(3) upgrades can be tested on Node~B before production sees them. A single
node with two or more GPUs can also train and serve at the same time by
dedicating GPUs to each, but it shares one failure domain and one power supply.

\ReportSubsection[sec:archive]{Archive knowledge: retrieval and fine-tuning}

Archived project data on the NAS should be usable by the model. Two techniques do
different jobs, and the design uses both.

\TableTitle{Retrieval versus fine-tuning for archive knowledge}
{\footnotesize
\begin{xltabular}{\textwidth}{@{}Y{0.632} Y{1.158} Y{1.053} Y{1.157}@{}}
  \toprule
  Technique & What it does & Strengths & Limits \\
  \midrule
  \endhead
  Retrieval (index the documents) & Splits documents into passages, stores them in a searchable index, and gives the model the best-matching passages with each question. & Quotes exact wording and cites document, clause and page; new files are usable as soon as they are indexed; follows NAS permissions; no training time. & Answer quality depends on finding the right passages; limited by the context size per question. \\
  Fine-tuning (train an adapter) & Adjusts the model using example questions and answers or documents. & Learns vocabulary, formats and how to do tasks (for example draft an inspection checklist in the house format); shorter prompts. & Does not reliably memorize a document and can state wrong clause numbers with confidence; content is hard to update or remove; must be re-tested after each change; the resulting weights are CUI. \\
  \bottomrule
\end{xltabular}}

Giving the model a PDF and training on it will not by itself make it a dependable
expert. For a standard such as NFPA~780, where exact wording matters, \textbf{retrieval
with citations is the primary mechanism}; fine-tuning on reviewed examples teaches the
model how to use and reason about the material. A fine-tuned adapter is deployed only
if it measurably beats retrieval alone.

\textbf{Example workflow: NFPA~780.}
\begin{enumerate}
  \item Place the licensed PDF in a controlled library folder on the NAS with the intended
    access groups.
  \item An ingestion job (run on Node~B or on CPU) extracts text and tables, applies optical
    character recognition to scanned pages, splits by section and clause, creates
    embeddings, and stores them in the vector database with edition and access metadata.
  \item Users ask questions; the model answers with citations to clause and page.
  \item Optional training: example questions, answers and tasks are generated from the
    document and reviewed by an engineer; an adapter is trained on Node~B; it is evaluated
    against a held-out set of questions with known answers and compared with retrieval alone.
  \item When a new edition is issued, re-ingest it with the edition recorded so that old and
    new requirements are not mixed.
\end{enumerate}

\TodoAnchor{aoa-nfpa}
\textbf{Points to resolve.}
\begin{itemize}
  \item \textbf{Formats.} Text-based and scanned documents work. CAD and drawing files (for
    example DWG or Revit) are not text; in the base plan only titles, annotations and
    metadata can be extracted. Reading drawings visually needs a vision model that is not
    budgeted here.
  \item \textbf{Publisher licenses.} Standards such as NFPA~780 are copyrighted. Confirm that
    the license allows internal indexing and AI training.
  \item \textbf{Access control.} Retrieval must respect NAS permissions. A model trained on a
    restricted folder must not be available to people without access, so train separate
    adapters per access group, or train only on material open to all users of that model.
  \item \textbf{Derived data.} Indexes, embeddings and adapters built from CUI are treated as
    CUI and covered by backup and sanitization rules.
  \item \textbf{Effort.} Ingestion, OCR and training run on Node~B and store results on
    Computer~1; archive size and formats (open item) set the effort.
\end{itemize}

\ReportSubsection[sec:nodeb]{Node B operating modes}
\TodoAnchor{aoa-nodeb}

Training alongside serving is the main use of Node~B. Because it is the same class of
machine as Node~A, it can also run a copy of the serving model. This is an option within
the design, not a fixed decision.

\TableTitle{Node B modes}
{\footnotesize
\begin{xltabular}{\textwidth}{@{}Y{0.656} Y{1.031} Y{1.313}@{}}
  \toprule
  Mode & When & Effect on capacity \\
  \midrule
  \endhead
  1. Training (default) & Fine-tuning, document ingestion, OCR and evaluation jobs run on Node~B. & Node~A serves users alone, so capacity is Node~A's. \\
  2. Second serving node & No training job is running; Node~B loads the same model and the gateway shares the load. & Adds about 3--8 simultaneous users (1 GPU) or 8--20 (2 GPUs). Training jobs take priority: the gateway drains Node~B sessions before a job starts. \\
  3. Failover & Node~A is unavailable. & Node~B serves alone at its own capacity (3--8 users with 1 GPU); training is paused. \\
  \bottomrule
\end{xltabular}}

A 1-GPU Node~B cannot train and serve at the same time: one 96~GB GPU holds either the
serving model with its cache (about 63~GB plus cache) or a training job (about 20--90~GB).
A 2-GPU Node~B (D4) can do one of each. Training jobs should save checkpoints often so
they can be paused for failover or for serving. A job queue (for example Slurm or
Kubernetes with priority classes) enforces the policy; a typical policy is training at
night and on weekends, and serving on Node~B during working hours when it is idle.

\ReportSubsection{Software stack (self-hosted, no outbound connectivity)}

\TableTitle{Software layers}
{\footnotesize
\begin{xltabular}{\textwidth}{@{}l Y{1.231} Y{0.769}@{}}
  \toprule
  Layer & Candidates & Purpose \\
  \midrule
  \endhead
  OS and virtualization & RHEL, Rocky or Ubuntu Pro in FIPS mode with STIG hardening; KVM/Proxmox; Docker or k3s & Hardened base; private cloud. \\
  Storage & ZFS (RAID-Z2) or equivalent, SMB/NFS, snapshot and replication tooling & File share, backups, datasets. \\
  Inference engine & vLLM or SGLang (TensorRT-LLM optional) & Continuous batching, paged KV cache, prefix caching. \\
  Training tools & PEFT/LoRA frameworks, experiment tracking (e.g., MLflow) & Adapter fine-tuning, evaluation. \\
  Gateway & LiteLLM or NGINX with OIDC, health checks, optional keepalived virtual IP & OpenAI-compatible API, quotas, failover routing. \\
  User interface & Open WebUI or similar & Chat, uploads, shared prompts. \\
  Retrieval & Open embedding model and reranker; pgvector or Qdrant & Cited answers with access-control enforcement. \\
  Document ingestion & PDF and Office text extraction, OCR for scanned pages, embedding model & Turn NAS archive content into searchable passages. \\
  Observability & Prometheus, Grafana, DCGM exporter, log forwarding to SIEM & Utilization, latency, audit. \\
  Configuration management & Ansible or similar, stored in the NAS git server & Identical node builds; fast rebuild. \\
  \bottomrule
\end{xltabular}}

\ReportSubsection{Model selection approach}
Use a portfolio sized to the node:

\TableTitle{Model roles and memory}
{\small
\begin{xltabular}{\textwidth}{@{}Y{0.844}Y{1.5}Y{0.656}@{}}
  \toprule
  Role & Candidate classes (verify versions, licenses, provenance) & Weights memory \\
  \midrule
  \endhead
  General assistant & OpenAI gpt-oss-120b (mixture of experts, Apache 2.0, about 5B active parameters) & about 63~GB (4-bit) \\
  Mid-size general or coding & Llama 3.x 70B, Gemma 27B, Mistral Small/Medium class, NVIDIA Nemotron & 15--40~GB (4-bit) to 70~GB (FP8) \\
  Large mixture of experts (4-GPU node only) & Llama 4 Maverick-class (about 400B total) at 4-bit & about 200--220~GB \\
  Embeddings and reranker & Small open models & 1--4~GB \\
  \bottomrule
\end{xltabular}}

\textbf{Provenance policy.} This report assumes U.S. or allied-origin models
only; many organizations prohibit models from foreign-adversary developers on
government-related systems regardless of license. Download on an internet-connected
system, scan, hash-verify, store as \texttt{safetensors} (no pickle files), and
import through the approved transfer process. Legal review of each license is
required. Mixture-of-experts models are well suited to this budget: few active
parameters per token mean one 96~GB GPU can serve several users quickly, and
the cache memory per token is small.

\ReportSubsection{How multiple users share the system}
\begin{enumerate}
  \item \textbf{Continuous batching:} the inference engine interleaves many users'
    requests token by token, so one loaded model serves everyone at once.
  \item \textbf{Paged KV cache:} each active conversation uses GPU memory in
    proportion to its context length; memory is usually the limit for long
    context (Section~\ref{sec:sizing}).
  \item \textbf{Prefix caching:} shared system prompts and repeated documents are cached.
  \item \textbf{Quotas and priority:} the gateway enforces per-user and per-team token
    budgets and prioritizes interactive use over batch jobs.
  \item \textbf{Isolation:} conversations and retrieval indexes are partitioned by
    user and group; logging is per user.
  \item \textbf{Batch queue:} large offline jobs run off-peak.
\end{enumerate}

\ReportSubsection[sec:resil]{Resilience, backup and multi-node growth}

\TableTitle{Failure modes and responses}
{\footnotesize
\begin{xltabular}{\textwidth}{@{}Y{0.8}Y{1.0}Y{1.0}Y{1.2}@{}}
  \toprule
  Failure & Single-node configs & Dual-node configs & Mitigation \\
  \midrule
  \endhead
  ML node A down & LLM unavailable. & Gateway fails over to Node~B; training paused. & Health checks; models staged locally on both nodes; documented switchover. \\
  ML node B down & Not applicable. & No user impact; no failover or training until repaired. & Spare parts; warranty. \\
  One GPU fails & Node continues with fewer GPUs if the model fits; otherwise down. & Same, plus failover. & Keep one cold-spare GPU if the budget allows. \\
  Computer 1 (NAS) down & File share, license servers, private cloud and the gateway/retrieval tier are down. & Same for NAS services; LLM serving continues if a second gateway runs on Node~B (retrieval corpus unavailable). & Redundant PSUs, RAID-Z2, UPS; VM backups; tested restore; optional replica (Section~\ref{sec:hw}). \\
  Disk failure on NAS & Array degraded but online. & Same. & RAID-Z2 tolerates two disks; hot spare or fast replacement. \\
  Power loss & All down. & All down. & UPS with graceful shutdown; generator if mission-critical. \\
  Bad model or software update & Rollback only. & Test on Node~B first. & Versioned registry; configuration as code. \\
  Site loss (fire, flood) & Data loss unless backed up elsewhere. & Same. & Encrypted offline or off-site copy of backups within the CUI boundary. \\
  \bottomrule
\end{xltabular}}

\textbf{Growth path (multi-node).}
\begin{itemize}
  \item \textbf{Same build, more nodes.} Every ML node uses the same chassis class,
    operating-system image and configuration, so a third or fourth node is another
    copy. The storage switch has spare 25~GbE ports, and the gateway load-balances
    across any number of nodes.
  \item \textbf{Scale out by replicas.} More nodes serve more simultaneous users
    (each node holds a full copy of the model) or run more training jobs in
    parallel. Add GPUs inside a node (about \$9K each) before adding a node.
  \item \textbf{Limits.} These PCIe-based workstation GPUs do not support
    efficient multi-node tensor parallelism, so a single model replica must fit
    inside one node (up to 384~GB on a 4-GPU node). A single training job spanning
    several nodes would need an RDMA-capable 100--200~GbE fabric (additional
    network cards and switch, roughly \$7--15K for two nodes; estimate); this is
    not part of the base plan.
  \item \textbf{Computer 1 redundancy.} A second NAS/private-cloud server (about
    \$12--18K) with replicated storage and virtual machines is the next step if
    the license servers or file share become critical; see Section~\ref{sec:open}.
\end{itemize}
